% Lift tijd v1.tex
% Stand-alone reproductie van een page-break / line-break probleem:
% de \textmodel-tabel binnen een subfigure (naast \graphicmodel) krijgt
% een onnodige afbreking als de figuur naar pagina 2 wordt geduwd door
% voorafgaande tekst en een question/solution-blok.
%
% Compileren: lualatex "Lift tijd v1.tex"

\documentclass{exam}
\printanswers
\usepackage[syntax=NL]{numodel}
\usepackage{subcaption}
\usepackage{geometry}
\geometry{a4paper,margin=2cm}

\numodelsetup{gridmaxx=4}

\begin{document}

\subsection*{Personenlift}

\newmodelprefix{lift}

\mvar{T}{t}{0}{\s}{2}{systeem}
\mvar{Dt}{dt}{0.03}{\s}{1}{systeem}
\mvar{Tmax}{T}{5.6}{\s}{2}{systeem}

\mvar{J}{j}{1.1}{\m\per\s\cubed}{2}{systeem}
\mvar{Da}{da}{\liftJ * \liftDt}{\m\per\s\squared}{2}{constante}
\mvar{G}{g}{9.81}{\m\per\s\squared}{3}{constante}
\mvar{Mcab}{m_{cab}}{700}{\kg}{3}{constante}
\mvar{Mcon}{m_{con}}{400}{\kg}{3}{constante}
\mvar{Frol}{F_{rol}}{40}{\N}{2}{constante}

\mvar{FzCab}{F_{Z,cab}}{}{\N}{3}{hulp}
\mvar{FzCon}{F_{Z,con}}{}{\N}{3}{hulp}
\mvar{FzNet}{F_{Z,net}}{}{\N}{3}{hulp}
\mvar{Fmotor}{F_{motor}}{}{\N}{3}{hulp}

\mvar{A}{a}{0}{\m\per\s\squared}{2}{voorraad}
\mvar{V}{v}{0}{\m\per\s}{1}{voorraad}
\mvar{H}{h}{0}{\m}{1}{voorraad}

\mrule{FzCab}{\liftMcab * \liftG}
\mrule{FzCon}{\liftMcon * \liftG}
\mrule{FzNet}{\liftFzCab - \liftFzCon}
\mrule*{A}{(\liftT < \liftTmax / 4) || (\liftT > \liftTmax * 3 / 4) ? \liftA + \liftDa : \liftA - \liftDa}
\mrule{V}{\liftV + \liftA * \liftDt}
\mrule[aliasright=\cdots]{H}{\liftH + \liftV * \liftDt}
\mrule[aliasright=\cdots]{Fmotor}{(\liftMcab + \liftMcon) * \liftA + \liftFzNet + \liftFrol}
\mrule{T}{\liftT + \liftDt}
\mstop{\liftT >= \liftTmax}

Een personenlift beweegt omhoog in een liftschacht. De lift hangt aan een
kabel die via een katrol verbonden is met een contragewicht. De katrol
wordt aangedreven door de liftmotor.

De gegevens over de lift staan in standaard eenheden bij de startwaarden
van het tekstuele model; $dt$ is de tijdstap waarmee het model rekent,
$v$ is de (start)snelheid, $a$ is de (start)versnelling, $h$ is de
(start)hoogte, $m_{\text{cab}}$ en $m_{\text{con}}$ zijn de massa's van
respectievelijk de cabine met passagiers en het contragewicht, $g$ is de
valversnelling, $T$ is de totale looptijd van het model, en
$F_{\text{rol}}$ is de constante wrijvingskracht in het liftsysteem. De
variabele $j$ noemt men de ruk (jerk), $j=a'(t)$ dat is de toename van
de versnelling per seconde, hieruit volgt de toename van de versnelling
per tijdstap, $da = j \cdot dt$.

Bij een systeem zonder aandrijving of rolwrijving bij de katrollen is de
spankracht hetzelfde in de hele kabel. In dit geval is er wel
aandrijving, daardoor kan er verschil zijn tussen de spankracht in de
kabel aan de kant van de cabine en aan de kant van het contragewicht.

Op een bepaald moment versnelt de lift met $\qty{0,3}{\m\per\second\squared}$
omhoog.

\begin{questions}

\question[4]
\textbf{Bereken} de spankracht aan de kant van het contragewicht.
\textbf{Noteer} de uitkomst met de juiste significantie.

\begin{solution}
  \begin{itemize}
    \item Inzicht dat $F_{res} = F_Z - F_{span}$
    \item Gebruik van $F_Z = m_{con}\,g$ en $F_{res} = m_{con}\,a$
    \item Completeren $F_{span} = F_Z - F_{res} = m_{con}\,g - m_{con}\,a$
    \item Significantie van twee (of consequent met alternatieve berekening).
  \end{itemize}
\end{solution}

De motor van de lift moet voldoende vermogen leveren om de cabine
omhoog te bewegen tegen de zwaartekracht in, en om de versnelling van
de cabine te realiseren.  Het contragewicht zorgt ervoor dat een
groot deel van het gewicht van de cabine wordt gecompenseerd, zodat
de motor minder kracht hoeft te leveren.  De netto kracht die de
motor moet leveren hangt af van het verschil tussen de
zwaartekrachten op cabine en contragewicht, de versnelling die de
lift maakt, en de rolwrijving in het systeem.

\question[2]
\textbf{Leg uit} waarom een contragewicht het energieverbruik van de
liftmotor verlaagt.

\begin{solution}
  \begin{itemize}
    \item Inzicht dat het contragewicht de zwaartekracht op de cabine
      gedeeltelijk compenseert.
    \item Completeren: de motor hoeft alleen de \emph{netto}
      zwaartekracht plus de versnellingskracht en de rolwrijving te
      leveren, niet de volle zwaartekracht op de cabine.
  \end{itemize}
\end{solution}

De snelheid en de hoogte van de lift als functie van de tijd worden
gemodelleerd met het computermodel in figuur~\ref{fig:lift-model}.

\begin{figure}[H]
\centering
\begin{subfigure}[c]{0.50\textwidth}
\centering
\textmodel[tblrenv=tblr]
\caption{Tekstueel model met SI-eenheden}
\label{fig:lift-model-tekstueel}
\end{subfigure}%
\hfill
\begin{subfigure}[c]{0.50\textwidth}
\centering
\graphicmodel
\caption{Grafisch model}
\label{fig:lift-model-grafisch}
\end{subfigure}
\caption{Computermodel voor de personenlift}
\label{fig:lift-model}
\end{figure}

\computemodel

\end{questions}

\end{document}
